\scanpage{181}
\mspage{494}
% position: first paragraph
Bien entendu, une telle valorisation ``officielle'' de l'action au détri-
ment de l'inaction, a aussitôt comme conséquence de mettre en branle en la per-
sonne des mécanismes de résistance (qui souvent restent occultes ou du moins
très brouillés), s'exprimant par une valorisation \mbox{o p p o s é e}: l'action\inserted{,} du
coup, apparaît comme ce qui est imposé par les dures nécessités de l'existence,
comme le \mbox{b o u l o t} en somme, emmerdant au possible, au bureau ou à l'usine
voire même aux champs, et éreintant en tous cas même s'il n'est pas trop emmer-
dant. La vraie raison d'être de l'action, c'est de gagner sa croûte et son logis
(ça c'est l'indispensable), et au-delà de ça et surtout, d'avoir des loisirs
chouettes (pendant sa vie active), et une retraite coquette et d'agréables loi-
sirs permanents plus tard, quand on sera dispensé de la regrettable obligation
``boulot''. Cette fois, c'est l'inaction (alias ``loisirs'') qui est valorisée de
façon plus ou moins consciente, et \inserted{c}'est l'action qui en est l'humble servante.
Il y a donc un \mbox{r e n v e r s e m e n t\quad d e s\quad r ô l e s}, mais avec tou-
jours le même déséquilibre: celui qui consiste en \mbox{l' a n t a g o n i s m e}
établi par l'intéressé (sous la poussée de conditionnements culturels) entre
deux aspects ou pôles essentiels de sa vie; antagonisme qui s'exprime et se
perpétue par un état de prépondérance despotique de l'un de ses aspects, et de
servitude de l'autre.

% position: second paragraph
Il me semble que le plus souvent, les deux attitudes et valorisations se
superposent chez une même personne, l'une dominant le pavé au niveau conscient,
l'autre au niveau inconscient. De la superposition de ces deux déséquilibres
opposés, visiblement, ne naît pas l'équilibre! Celui-ci par contre découle na-
turellement d'une compréhension de la vraie nature de l'action et de l'inaction
(même lorsqu'une telle compréhension reste purement ``instinctive'', se manifes-
tant directement par un comportement équilibré, et nullement par un ``savoir''
verbalisé). \mbox{D a n s\quad l' a c t i o n\quad a u\quad p l e i n\quad s e n s\quad d u\quad t e r-}
\mbox{m e}, il \mbox{y\quad a\quad a u s s i\quad l' i n a c t i o n} - elle y est \mbox{d a n s}
\mbox{l' i n s t a n t\quad m ê m e} j'entends, et pas seulement ``après'', parce qu'il
faut bien se reposer après l'action! Cette ``inaction'' dans l'``action'', le ``yin
dans le yang'' donc, est comme un calme profond qui sert d'assise à un mouvement
qui aurait lieu à la surface. Elle se manifeste par exemple par l'impression
de détente parfaite qui se dégage d'un félin en mouvement, que ce soit le pre-
mier chat de gouttière venu, ou une lionne à la carrure puissante...

% position: third paragraph
Et de même \mbox{d a n s\quad l' i n a c t i o n\quad v é r i t a b l e}, fût-elle
même totale, \mbox{i l\quad y\quad a\quad a c t i o n}. Ainsi le \mbox{s o m m e i l} est riche
de ses rêves qui nous parlent de nous-même, par lesquels nous vivons une
% position: foot pencil
\margin{180}
